%% The following is a directive for TeXShop to indicate the main file
%%!TEX root = main.tex

\chapter{Introduction}
\label{ch:introduction}

Replace this text with your introduction.  A typical introduction
motivates the problem area, identifies the gap in existing work,
states the thesis statement, and previews the contributions of each
research chapter~\citeeg{knuth1984literate}.

\section{Motivation}

Describe the problem your dissertation addresses and why it matters.
Cite prior work with \verb|\citep{}| for parenthetical citations
\citep{lamport1994latex} or \verb|\citet{}| for textual ones.

\section{Thesis Statement and Research Questions}

State your thesis in one or two sentences.  Then enumerate the
high-level research questions.  The \env{researchquestions}
environment (defined in \file{macros.tex}) numbers them globally:

\begin{researchquestions}
\item How can the first aspect of the problem be addressed?
\item How effective is the proposed approach compared to the state of
the art?
\end{researchquestions}

\section{Contributions and Organization}

Summarize the contribution of each chapter and how the chapters fit
together.  \autoref{chap:example-one} presents the first
contribution.  \autoref{chap:example-two} presents the second
contribution.  \autoref{ch:discussion} concludes the dissertation and
discusses future directions.
